% ═══════════════════════════════════════════════════════════════
% PREAMBLE — ALL PACKAGES & GLOBAL CONFIGURATION FOR TCC 1 UFN
% ═══════════════════════════════════════════════════════════════

% ── ENCODING & LANGUAGE ──────────────────────────────────────
\usepackage[utf8]{inputenc}
% UTF-8 encoding para caracteres em português (ã, ç, ê, etc.)

\usepackage[T1]{fontenc}
% T1 font encoding para hyphenação correta de palavras acentuadas

\usepackage[brazilian]{babel}
% Regras de hyphenação português brasileiro e formatação de data

% ── TYPOGRAPHY ───────────────────────────────────────────────
\usepackage{mathptmx}
% Times New Roman para texto e matemática — padrão UFN

% ── GEOMETRY (MARGINS UFN) ───────────────────────────────────
\usepackage[
  a4paper,
  top=2.5cm,
  bottom=2.5cm,
  left=1.8cm,
  right=1.8cm
]{geometry}
% Sobrescreve defaults do IEEEtran para atender margens UFN

% ── FIGURES ──────────────────────────────────────────────────
\usepackage{graphicx}
% \includegraphics — figuras individuais da pasta /images

\usepackage{subcaption}
% \begin{subfigure} — figuras lado a lado (ex.: múltiplas views)

% ── MATHEMATICS ──────────────────────────────────────────────
\usepackage{amsmath}
% align, equation, matrix environments — para fórmulas IA/IK

\usepackage{amssymb}
% \mathbb{R}, \forall, \exists — para notação formal

% ── CITATIONS ────────────────────────────────────────────────
\usepackage{cite}
% Estilo IEEE com compressão numérica: [1]–[4] ao invés de [1][2][3][4]

% ── CODE LISTINGS ────────────────────────────────────────────
\usepackage{listings}
\usepackage{xcolor}
% Configure listings para código monoespacial com números de linha:
\lstset{
  basicstyle=\ttfamily\footnotesize,
  numbers=left,
  numberstyle=\tiny\color{gray},
  stepnumber=1,
  numbersep=5pt,
  backgroundcolor=\color{gray!10},
  showspaces=false,
  showstringspaces=false,
  showtabs=false,
  frame=single,
  rulecolor=\color{black!30},
  tabsize=2,
  captionpos=b,
  breaklines=true,
  breakatwhitespace=false,
  keywordstyle=\color{blue!70}\bfseries,
  commentstyle=\color{green!50!black}\itshape,
  stringstyle=\color{red!60},
  language=Python,
  literate=
    {á}{{\'a}}1 {é}{{\'e}}1 {í}{{\'i}}1 {ó}{{\'o}}1 {ú}{{\'u}}1
    {Á}{{\'A}}1 {É}{{\'E}}1 {Í}{{\'I}}1 {Ó}{{\'O}}1 {Ú}{{\'U}}1
    {ã}{{\~a}}1 {õ}{{\~o}}1 {ç}{{\c{c}}}1 {ñ}{{\~n}}1
}
% Bloco literate trata caracteres portugueses dentro de blocos de código

% ── TABLES ───────────────────────────────────────────────────
\usepackage{tabularx}
% Colunas auto-expansíveis — necessário para tabela de cronograma TCC 2

\usepackage{booktabs}
% \toprule, \midrule, \bottomrule — regras profissionais (sem \hline)

% ── PDF METADATA & HYPERLINKS ────────────────────────────────
\usepackage[
  pdftitle={TCC 1 — [TÍTULO DO TRABALHO]},
  pdfauthor={[[NOME DO ALUNO]}],
  pdfsubject={Trabalho de Conclusão de Curso — UFN},
  pdfkeywords={[palavra1], [palavra2], [palavra3]},
  colorlinks=true,
  linkcolor=black,
  citecolor=black,
  urlcolor=blue!70!black,
  bookmarks=true
]{hyperref}
% DEVE ser o ÚLTIMO pacote carregado — trata metadados PDF e refs clicáveis

% ── CUSTOM ABNT COMMANDS ─────────────────────────────────────
% ABNT NBR 10520 — Citação longa (> 3 linhas):
% 4cm recuo esquerdo, fonte \small, espaçamento simples, sem aspas
\newenvironment{citacaolonga}{%
  \begin{list}{}{%
    \setlength{\leftmargin}{4cm}%
    \setlength{\rightmargin}{0cm}%
    \setlength{\topsep}{0pt}%
    \setlength{\partopsep}{0pt}%
  }%
  \item\small\setstretch{1.0}%
}{%
  \end{list}%
}
% Uso: \begin{citacaolonga} texto da citação... \end{citacaolonga}

% Macro de conveniência para nome institucional
\newcommand{\UFN}{Universidade Franciscana (UFN)}